\section{Research target, source state, and contribution}
\label{sec:scope}

\subsection{The particular obstruction addressed here}
The maintained algorithm has several mathematically distinct layers.
The complete recognizer can use exact invariant computations and other
certifying procedures.  The normal-surface research pipeline separately
constructs compressed cuts, their parallelity bundles, finite complementary
cores, ordered interfaces, and, when a sufficient criterion succeeds,
handlebody censuses.  A census records every component's genus and binary
multiplicity.  The latter pipeline supplies potential building blocks for
hierarchies, but it is not itself the complete unknot recognizer.

This work fixes one precise gap in that pipeline.  At the source revision
\basecommit\ \cite{ProveIt}, both the core certifier and the supplied-cut entry point
stop on a nonempty list \code{handles3}.  This is a conservative implementation
restriction, not a theorem that the corresponding manifolds are difficult.
The retained source charts already contain sufficient information to
remove these handles geometrically.  Our first contribution is to recover
that information, prove the reduction relative to the exposed interfaces,
and connect it to the existing independently checked boundary kernel.

The pinned files are in
\code{Topology/UnknotRecognition/fast/fastunknot/}; the principal entry point is
\code{normal_cut_handlebodies.py}.  The new branch reuses its already computed
parallelity packet.  It does not start a second parallelity discovery.
The new certificate is dispatched through
\code{normal-cut-handlebodies-v6}.  Earlier schemas remain supported.
The package includes the patch and a machine-readable source manifest.

\subsection{Why this is a continuation of the existing work}
We inspected the current implementation, the maintained synthesis, and the
ten archives in \code{docs/incoming} at the pinned revision.  Several tempting
ideas are already present there and are not claimed as discoveries here.
In particular:
\begin{itemize}
\item unsigned stopping sets and complete unit-column peeling are established
  by the incoming stopping-core and saturated-bundle reports;
\item cographic persistence uses chamber graphs to retain simultaneously
  nonseparating surfaces and accounts for causal resets;
\item modular Morse certificates verify supplied compressed matchings, with
  an explicit regular-geometry hypothesis;
\item genus budgets, sparse cut grammars, and native ordered interfaces already
  give several complementary size and attachment bounds.
\end{itemize}
None of those ten packages implements the native dual two-/three-handle
normalization developed below.  The combination of dual forest selection
with the stopping obstruction also poses a new optimization problem within
this development.  Its degree-two-column case leads to the polynomial
algorithm in Section~\ref{sec:sparse}.

Our novelty claims are deliberately local: a source-specific compiled
normalization and certificate interface, together with the stated
constrained-cancellation characterizations and algorithms.  We make no claim
that handle cancellation, spanning trees, or matroid intersection are new,
and no claim of an exhaustive priority search across all literature.

\subsection{Relationship to Lackenby's programme}
Lackenby's 2026 preprint \cite{Lackenby2026} presents a hierarchy algorithm
for essential boundary patterns and incompressibility.  Its
Proposition~9.1 bounds iterations by
\begin{equation}\label{eq:lackenby-iterations}
                 L(g+1)^L,
\end{equation}
where $L$ bounds hierarchy length and $g$ bounds surface pattern-complexity.
Its Section~9 explicitly separates this count from a full running-time
analysis.  The refined construction gives a linear bound on hierarchy
length in Proposition~10.2; the introduction leaves the intended speed-up
for further work.

There is also a distinction between two announced exponents.  Early Oxford
publicity uses $n^{O(\log n)}$ \cite{Oxford2021}.  The author's March 2021 slides
give $L=O((\log n)^2)$ and $g=O(n^2)$, resulting in
\begin{equation}\label{eq:announced}
       n^{O((\log n)^2)}=2^{O((\log n)^3)};
\end{equation}
the 2024 Tsinghua seminar description repeats this latter bound
\cite{LackenbySlides,Tsinghua2024}.  Both are quasi-polynomial, but they
are not identical.  We aim toward the general quasi-polynomial programme
without treating a local polynomial operator as a proof of either exponent.

Compressed cutting itself is established prior art.  Theorem~9.3 of
\cite{Lackenby2021} already gives a polynomial algorithm in triangulation
size and logarithmic surface weight, using Agol--Hass--Thurston orbit
methods \cite{AHT2006}.  Theorem~14.4 of \cite{Lackenby2026} gives the
uniform-handle version, including bundle type and attachment information.
Our contribution concerns an explicit subsequent native operator and its
certification.  It does not replace or re-prove those general existence
theorems.

\subsection{Statements proved and statements left open}
The main mathematical results are as follows.
\begin{enumerate}
\item The native dual graph has $O(t)$ vertices, edges, and incidence work.
  It is rooted for every admitted source and every admitted normal vector.
  Its spanning trees give geometric reductions of all core three-handles
  (Theorem~\ref{thm:native}).
\item This graph depends only on which source edge weights vanish.  A tree
  can therefore be reused throughout such a support stratum, including all
  positive scalar multiples (Corollary~\ref{cor:support}).
\item For a fixed lower spanning forest, successful dual trees are exactly
  the graphic bases that intersect every stopping set of the remaining
  unsigned occurrence matrix (Theorem~\ref{thm:transversal}).
\item Stopping witnesses give a complete parameterized search; insufficient
  graphic rank of a family of disjoint witnesses gives a rigorous obstruction
  (Theorem~\ref{thm:search} and Proposition~\ref{prop:packing}).
\item If every non-inert column has at most two incident rows, tree selection
  is polynomial and has explicit positive and negative certificates
  (Theorem~\ref{thm:degree-two}).  Every original native column has at most
  three total occurrences, isolating $(1,1,1)$ as the remaining native
  type (Proposition~\ref{prop:three}); its odd entries form a dual cycle
  (Proposition~\ref{prop:cycle}).
\end{enumerate}
The geometric scope is a supplied compressed cut with authenticated charts.
The exact search scope fixes the lower zero-/one-handle forest, allows the
specified dual-tree reduction, and then allows raw unsigned one-/two-handle
cancellation.  Failure of this restricted search proves no assertion that
the manifold is not a handlebody or that the original knot is nontrivial.
Arbitrary handle slides, changes of the lower forest, and hierarchy-wide
search remain outside that completeness theorem.
